% Direct-authority diplomatic French for printed p.159.
\SourcePage{159}

\LayerTarget{tohoku:def:flasque-and-soft-sheaves}
Un faisceau d'ensembles $F$ sur un espace $X$ est dit \textit{flasque}
(resp. \textit{mou}) si pour toute partie $A$ ouverte (resp. fermée) de $X$,
toute section de $F$ sur $A$ est la restriction d'une section de $F$ sur
$X$. Ainsi si on se donne pour tout $x\in X$ un ensemble $E_x$, le faisceau
$E$ dont l'ensemble des sections sur un ouvert $U$ est
\[
E(U)=\prod_{x\in U}E_x
\]
(avec les applications de restrictions évidentes ; comparer l'exemple traité
avant prop. 3.1.2) est évidemment flasque et mou ; on en conclut que tout
faisceau d'ensembles se plonge dans un faisceau flasque d'ensembles, de même
tout faisceau de groupes abéliens (resp. de $\mathbf O$-modules, si
$\mathbf O$ est un faisceau d'anneaux sur $X$) se plonge dans un faisceau
flasque de groupes abéliens (resp. de $\mathbf O$-modules). Notons que si une
partie fermée $A$ de $X$ admet un voisinage paracompact, alors toute section
sur $A$ d'un faisceau $F$ défini sur $X$ est la restriction d'une section
définie dans un voisinage de $A$ ; il en résulte aussitôt que
\textit{si $X$ est paracompact, un faisceau flasque est mou}.

\LayerTarget{tohoku:def:phi-soft-and-paracompactifying}
Soit maintenant $\Phi$ une famille de parties fermées de $X$ satisfaisant aux
conditions générales envisagées dans 3.2, un faisceau $F$ de groupes abéliens
sur $X$ est dit $\Phi$-\textit{mou} si pour tout $A\in\Phi$ et toute section
de $F$ sur $A$, il existe une $f\in\Gamma_\Phi(F)$ induisant la section
donnée. Disons que $\Phi$ est une \textit{famille paracompactifiante} si, en
plus des conditions déjà envisagées, elle satisfait aux conditions
supplémentaires suivantes (introduites d'abord dans [4]) :
\textit{tout $A\in\Phi$ est paracompact, et admet un voisinage $B\in\Phi$}.
On voit alors facilement, comme ci-dessus, \textit{qu'un faisceau flasque de
groupes abéliens est $\Phi$-mou pour toute famille $\Phi$
paracompactifiante}, d'où résulte en particulier que tout faisceau de groupes
abéliens se plonge dans un faisceau $\Phi$-mou (puisqu'il se plonge même dans
un flasque). L'intérêt des définitions précédentes tient à la

\LayerTarget{tohoku:prop:3.3.2}
\textsc{Proposition 3.3.2.} --- \textit{Soit $X$ un espace topologique muni
d'une \og famille $\Phi$\fg{}. Considérons le foncteur $\Gamma_\Phi$ défini
sur la catégorie $C^X$ des faisceaux de groupes abéliens sur $X$, et à
valeurs dans la catégorie des groupes abéliens. Alors les conditions du lemme
3.3.1 sont satisfaites dans chacun des deux cas suivants :}

\textit{1) $\mathbf M$ est la famille des faisceaux flasques sur $X$.}

\textit{2) $\Phi$ est paracompactifiante, et $\mathbf M$ est la famille des
faisceaux $\Phi$-mous sur $X$.}

\LayerTarget{tohoku:cor:3.3.2}
\textsc{Corollaire.} --- \textit{$H^p_\Phi(X,F)=0$ pour $p>0$ si $F$ est
flasque, ou si $F$ est $\Phi$-mou et la famille $\Phi$
paracompactifiante.}

La condition (i) du lemme a déjà été vérifiée, la condition (ii) se vérifie
trivialement, seule la condition (iii) demande une démonstration, pour
laquelle nous renvoyons au livre de Godement (ou que nous proposons au
lecteur à titre d'exercice).

\LayerTarget{tohoku:remark:fine-sheaves-are-soft}
\textsc{Remarque.} --- Si la famille $\Phi$ est paracompactifiante, on
vérifie facilement que les faisceaux fins [4, XV] sont $\Phi$-mous, donc
vérifient $H^p_\Phi(X,F)=0$. Il en résulte que la théorie de la cohomologie
donnée dans [4] pour des familles $\Phi$ paracompactifiantes est bien un cas
particulier de celle développée ici. Nous renvoyons aussi au livre de Godement
pour une définition, particu-

\TohokuContinuation{160}

